\section{Discussion}\label{sec:discussion}

\paragraph{What the pieces say together.} The baseline shows what independent corroboration is worth when labels are accurate. RQ1 shows that this value is not a property of labelling in general but of one number, $J=r-\varphi$: it turns negative when labels report false sharing more often than they detect real sharing. RQ2 shows that real metadata is mostly \emph{absent} and, where present, mostly correct. RQ3 shows that LLMs do use a label that is present, but in the direction their instruction dictates.

\paragraph{Answers.}
\emph{RQ1:} metadata helps if and only if $J>0$ and harms when $J<0$; this is proved within the model and confirmed in simulation on new tasks, and missing labels shrink the effect without flipping it. \emph{RQ2:} among real MCP servers, provider sharing is real (16--27\% of pairs) and varies strongly across samples; metadata is missing in about 70\% of pairs; when present it is mostly right ($J>0$), but usually only at the operator level. \emph{RQ3:} with a prompt that asks for an independent source and complete, correct labels, two small LLMs choose a shared-provider corroborator far less often ($0.42\to0.11$ and $0.34\to0.05$), and a larger model agrees; with a neutral prompt the same label raises that rate to about $0.80$. The inversion experiment shows that the models act on the label's value, avoiding the candidate that carries the primary's provider when asked for independence and seeking it when asked for consistency.

\paragraph{Implications} (inferences that need further validation).
\begin{enumerate}
  \item \textbf{The bottleneck is missing metadata, not wrong metadata.} Because $J>0$, real metadata does not harm selection under the theorem; it is simply too sparse. Moving the provider name from the README into the tool description is a cheap change that, in the projection, lowers the chance of choosing a shared corroborator from 0.289 to 0.109. This is a concrete recommendation for MCP server authors and tool registries. RQ3 partly checks it: with complete labels and a prompt asking for independence, the shared-choice rate of the models falls to $0.04$--$0.11$.
  \item \textbf{A label helps only if the agent is told to use it for independence.} A neutral request to ``double-check'' makes the models read the same label as a sign of a consistent tool. Registries that publish providers should come with explicit guidance, in agent instructions, to look for a \emph{different} source.
  \item \textbf{Operator names are not enough.} At the data-origin level, provider names miss about 17\% (survey 2, preregistered) to 39\% (survey 1, post hoc) of dependencies; Frankfurter distributes ECB data, and differently named web-search servers call the same engines. Trustworthy registries should also publish data origins.
  \item \textbf{Independent sources can be scarce.} In web search, every surveyed server shares an engine with at least one other server. Good metadata then helps only partially, and an agent should recognise when no independent source exists and escalate, rather than treating agreement between two tools as strong evidence.
\end{enumerate}
